% ============================================================
\chapter{練習問題と答え}\label{app:exercises}

本文の式を手で使ってみるための問題である。答えは問題のあとにまとめた。電卓があれば解ける。

\section{問題}

\begin{enumerate}[label=\textbf{問 \arabic*.},leftmargin=3.2em]
\item \textbf{採用率}：ある株式の年率の標準偏差は 20\%、$\tau=0.05$ とする。この株式に絶対ビューを 1 つ置き、ビューの不確かさを標準偏差で 4\%（$\omega=0.04^2$）とした。採用率はいくつか。

\item \textbf{Idzorek の閉じた式}：問 1 と同じ株式で、確信度 $C=20\%$ にしたい。$\omega$ をいくつにすればよいか。そのときの採用率が 20\% になることも確かめよ。

\item \textbf{IC から $\omega$}：月次の標準偏差が 6\% の銘柄について、IC が 0.08 の予想がある。市場版の BL で、当たり具合に見合う $\omega$ と採用率はいくつか。点数が平均より 1.5 標準偏差高い銘柄の、期待リターンの上乗せ（Grinold の式）はいくつか。

\item \textbf{$\tau$ が消えること}：He--Litterman の $\Omega=\tau\diag(P\Sigma P\T)$ で、$\tau$ を 0.025 から 0.1 に変えた。事後の期待リターンと、事後の不確かさ $\sqrt M$ はそれぞれどう変わるか。

\item \textbf{確信の積み上がり}：200 銘柄に絶対ビューを置く。リターンの相関はどの組でも 0.25、1 銘柄だけなら採用率 1\% になる $\omega$（$\omega=99\sigma^2$）を、対角の $\Omega$ で使う。全体の水準（平均の方向）の採用率はいくつか。外れの相関も 0.25 だと正しく入れたら、いくつになるか。

\item \textbf{較正の傾き}：予想と翌月のリターンの相関が 0.05、翌月のリターンの横断面の標準偏差が 8\%、予想の標準偏差が 2\% だった。較正の傾き $b$ はいくつか。予想をそのまま期待リターンにすると、何倍の自信過剰か。

\item \textbf{$\tau\approx1/T$}：月次の共分散を使い、均衡リターンを「5 年分の月次データの平均と同じくらい確か」とみなすなら、$\tau$ はいくつか。

\item \textbf{TE と IR}：能動運用の最適化で、最大の情報レシオが 0.4 だった。TE の目標を年率 2\% にすると、期待超過リターンと、TE への回避度 $\lambda_A$ はいくつか。

\item \textbf{エントロピー・プーリング}：2 つのシナリオ（$-10\%$ と $+10\%$）が確率 0.5 ずつの事前分布に、「期待リターンは $+2\%$」というビューを入れる。事後の確率と、相対エントロピーはいくつか。

\item \textbf{CVaR}：等しい確率の 20 個のシナリオの損失を小さい順に並べると、いちばん大きいものが 10\%、2 番目が 6\% だった。95\% の VaR と CVaR はいくつか。

\item \textbf{ESG の共分散の案}：ESG の点数が 0.25 の銘柄に、式~\eqref{eq:esg-penalty} を当てはめると、分散と標準偏差はそれぞれ何倍になるか。均衡も同じ共分散で作り直したら、ビューがないときの比率はどうなるか。

\item \textbf{モメンタムのボラティリティ調整}：前月までの 126 日の日次リターンから見積もったモメンタムの月次の標準偏差が 8\% だった。目標を年率 12\% とすると、倍率はいくつか。
\end{enumerate}

\section{答え}

\begin{enumerate}[label=\textbf{問 \arabic*.},leftmargin=3.2em]
\item $\tau\sigma^2=0.05\times0.04=0.002$、$\omega=0.0016$。採用率は $0.002/(0.002+0.0016)=55.6\%$。ビューのほうを少し強く信じる。

\item $\omega=\tau\sigma^2(1-C)/C=0.002\times0.8/0.2=0.008$。採用率は $0.002/(0.002+0.008)=20\%$ で、確信度と一致する（\secref{sec:idz-closed}）。

\item $\omega=\sigma^2(1/\IC^2-1)=0.0036\times(156.25-1)\approx0.559$。採用率は $\IC^2=0.64\%$。上乗せは $\IC\cdot\sigma\cdot z=0.08\times6\%\times1.5=0.72\%$（月次）。$\omega$ の平方根は約 75\% で、「来月のリターンの予想は ±75\% 外れうる」と言っていることになる。それくらい雑音の多い予想だ、ということである。

\item 事後の期待リターンは変わらない（$\tau$ が約分される）。$M$ は $\tau$ に比例するので、$\sqrt M$ は $\sqrt{0.1/0.025}=2$ 倍になる（\secref{sec:tau-cancel}）。

\item $\sigma^2=1$ とおくと、平均の方向の $\Sigma$ の固有値は $1-0.25+0.25\times200=50.75$。対角の $\Omega$ の固有値は 99 なので、採用率は $50.75/(50.75+99)\approx33.9\%$。外れの相関も 0.25 なら、$\Omega$ の固有値は $99\times50.75$ になり、採用率は $50.75/(50.75+99\times50.75)=1/100=1\%$ のまま（\chapref{ch:stacking}）。

\item $b=\IC\cdot\sigma_r/\sigma_q=0.05\times8/2=0.2$。そのまま使うと、1 ÷ 0.2 = 5 倍の自信過剰。

\item $T=60$ か月なので、$\tau\approx1/60\approx0.017$。

\item 期待超過リターンは $\IR\times\TE=0.4\times2\%=0.8\%$。$\lambda_A=\IR^*/\TE^*=0.4/0.02=20$（\chapref{ch:te}）。

\item 事後の確率を $(q_1,q_2)$ とすると、$-0.1q_1+0.1q_2=0.02$、$q_1+q_2=1$ から、$q=(0.4,0.6)$。相対エントロピーは $0.4\log(0.4/0.5)+0.6\log(0.6/0.5)=0.4\times(-0.223)+0.6\times0.182\approx0.020$。

\item 損失が 6\% 以下になる確率は $19/20=95\%$ なので、95\% の VaR は 6\%。CVaR は、Rockafellar と Uryasev の定義（損失の大きい側の確率 5\% の部分の平均）では、いちばん大きいシナリオ 1 個の 10\% である。「VaR 以上の損失の平均」と読むと、6\% と 10\% の平均の 8\% になる。離散のシナリオでは定義によって値が変わり、シナリオが少ないと 1〜2 個のシナリオで決まってしまう（\chapref{ch:cvar} の推定の誤差の節）。

\item 分散は $2-0.25=1.75$ 倍、標準偏差は $\sqrt{1.75}\approx1.32$ 倍。均衡も同じ共分散で作り直すと、ビューがないときの比率は時価総額の比率のままで、何も変わらない（\secref{sec:esg-penalty}）。

\item 目標の月次の標準偏差は $12\%/\sqrt{12}\approx3.46\%$。倍率は $3.46/8\approx0.43$ で、ポジション（ビューの確信度）を半分以下にする（\chapref{ch:momentum}）。
\end{enumerate}
